%----------------------------------------------------------------------
\subsection{Calculation of the vacuum incoming state}
\label{sec:wf-vac}
%----------------------------------------------------------------------

We now compute the LCWF components that fix the coefficients. Throughout, the outgoing quark is
$\qk{\beta}{\lambda_{1}}{l}$ and the outgoing antiquark $\aqk{\alpha}{\lambda_{2}}{m}$, the dashed
vertical lines denote the intermediate states, the double line denotes the valence charge, a
straight line crossed by a bar denotes an instantaneous exchange, and we abbreviate
$E_{g}(k)=\eps_{k}$, $E_{q\bar q}(l,m)=\eps_{l}+\eps_{m}$, $E_{gg}(k,p)=\eps_{k}+\eps_{p}$. The
matrix elements are collected in appendix~\ref{app:me}. Each wave function is labelled, as in
Paper~I, by its outgoing particles and its powers of $\rho$.

\subsubsection{Quark--antiquark outgoing}

The order $g$ component of the vacuum is the Weizs\"acker--Williams gluon of Paper~I. A soft pair
appears first at order $g^{2}$, through two diagrams.

\paragraph{(a) The valence emits a gluon which splits into a pair.}
\begin{center}
\begin{tikzpicture}
\draw[val] (0,0) -- (3.9,0); \node[left] at (-0.05,0) {\scriptsize$\rho$};
\draw[gl] (1.1,0) -- (2.1,0.85); \node[above left] at (1.75,0.55) {\scriptsize$\gl{a}{i}{k}$};
\draw[qk] (2.1,0.85) -- (3.9,1.6); \node[right] at (3.9,1.6) {\scriptsize$\qk{\beta}{\lambda_{1}}{l}$};
\draw[qk] (3.9,0.3) -- (2.1,0.85); \node[right] at (3.9,0.3) {\scriptsize$\aqk{\alpha}{\lambda_{2}}{m}$};
\draw[cut] (1.6,-0.35) -- (1.6,1.95); \draw[cut] (3.0,-0.35) -- (3.0,1.95);
\end{tikzpicture}
\end{center}
Applying the LCPT rules,
\begin{equation}
\ket{\Psi^{1}_{q\bar q\,\rho}}=\int\! d^{3}k\,d^{3}l\,d^{3}m\;
\ket{\qk{\beta}{\lambda_{1}}{l}\aqk{\alpha}{\lambda_{2}}{m}}\;
\frac{\bra{\qk{\beta}{\lambda_{1}}{l}\aqk{\alpha}{\lambda_{2}}{m}}H_{gq\bar q}\ket{\gl{a}{i}{k}}}
{-E_{q\bar q}(l,m)}\;
\frac{\bra{\gl{a}{i}{k}}H_{g}\vac}{-E_{g}(k)} ,
\label{eq:Pvac1def}
\end{equation}
where here and below the sums over colours, polarizations, helicities and flavours are implicit.
Using the matrix elements \eqref{eq:MEqqg} and \eqref{eq:MEg} and $k=l+m$,
\begin{equation}
\ket{\Psi^{1}_{q\bar q\,\rho}}=\frac{g^{2}}{32\pi^{3}}\int\!d^{3}l\,d^{3}m\;
\ket{\qk{\beta}{\lambda_{1}}{l}\aqk{\alpha}{\lambda_{2}}{m}}\;
\frac{t^{a}_{\beta\alpha}\,\rho^{a}(-k)\,k^{i}\,\Ph^{i}_{\lambda_{1}\lambda_{2}}(k;l,m)}
{(k^{+})^{2}\,\eps_{k}\,\big(\eps_{l}+\eps_{m}\big)} .
\label{eq:Pvac1}
\end{equation}

\paragraph{(b) Instantaneous pair creation off the valence charge.}
\begin{center}
\begin{tikzpicture}
\draw[val] (0,0) -- (3.9,0); \node[left] at (-0.05,0) {\scriptsize$\rho$};
\draw[ins] (1.8,0) -- (1.8,0.85); \ibar{1.8}{0.45}
\draw[qk] (1.8,0.85) -- (3.9,1.6); \node[right] at (3.9,1.6) {\scriptsize$\qk{\beta}{\lambda_{1}}{l}$};
\draw[qk] (3.9,0.3) -- (1.8,0.85); \node[right] at (3.9,0.3) {\scriptsize$\aqk{\alpha}{\lambda_{2}}{m}$};
\draw[cut] (3.0,-0.35) -- (3.0,1.95);
\end{tikzpicture}
\end{center}
This is the $\rho\cdot J_{q}$ term of $H_{q-{\rm inst}}$, Eq.~\eqref{eq:Hqinst}:
\begin{multline}
\ket{\Psi^{2}_{q\bar q\,\rho}}=\int\! d^{3}l\,d^{3}m\;
\ket{\qk{\beta}{\lambda_{1}}{l}\aqk{\alpha}{\lambda_{2}}{m}}\;
\frac{\bra{\qk{\beta}{\lambda_{1}}{l}\aqk{\alpha}{\lambda_{2}}{m}}H_{q-{\rm inst}}\vac}{-E_{q\bar q}(l,m)}
\\
=-\frac{g^{2}}{8\pi^{3}}\int\!d^{3}l\,d^{3}m\;
\ket{q\bar q}\;
\frac{t^{a}_{\beta\alpha}\,\rho^{a}(-k)\,\delta_{\lambda_{1}\lambda_{2}}}
{(k^{+})^{2}\,\big(\eps_{l}+\eps_{m}\big)} ,
\label{eq:Pvac2}
\end{multline}
with $k\equiv l+m$ and the matrix element \eqref{eq:MEinst}.

\subsubsection{Summary of results}
\label{sec:sum-vac}

Using
$k^{i}\Gam^{i}(k;l,m)=4\eps_{k}-2(\eps_{l}+\eps_{m})-\tfrac{k^{+}}{l^{+}m^{+}}(\bm\sigma\!\cdot\!\bm l)(\bm\sigma\!\cdot\!\bm m)$,
which follows from \eqref{eq:Phi} and $\bm k=\bm l+\bm m$, the term proportional to $\eps_{k}$ in
\eqref{eq:Pvac1} cancels the instantaneous contribution \eqref{eq:Pvac2} exactly, and the quark
component of the vacuum at order $g^{2}$ is
\begin{equation}
\ket{\Psi_{q\bar q\,\rho}}=-\frac{g^{2}}{8\pi^{3}}\int\!d^{3}l\,d^{3}m\;
\ket{\qk{\beta}{\lambda_{1}}{l}\aqk{\alpha}{\lambda_{2}}{m}}\;
t^{a}_{\beta\alpha}\rho^{a}(-k)\left[\frac{\delta_{\lambda_{1}\lambda_{2}}}{k^{+}\bm k^{2}}
+\frac{\chi^{\dg}_{\lambda_{1}}(\bm\sigma\!\cdot\!\bm l)(\bm\sigma\!\cdot\!\bm m)\chi_{\lambda_{2}}}
{2\,l^{+}m^{+}\bm k^{2}\,(\eps_{l}+\eps_{m})}\right].
\label{eq:Pvac}
\end{equation}
The cancellation of the instantaneous term is required by unitarity: the relation
\eqref{eq:UA2} contains no instantaneous piece. Since the pair component is of order $g^{2}$, its
norm is of order $g^{4}$, and the normalization $N$ of the vacuum is not modified by quarks at the
order we work. The full vacuum state through $\order{g^{2}}$ is the one of Paper~I supplemented by
\eqref{eq:Pvac}, with $\rho=\rho_{g}+\rho_{q}$ in all the gluonic components.

%----------------------------------------------------------------------
\subsection{Calculation of the single-gluon incoming state}
\label{sec:wf-g}
%----------------------------------------------------------------------

We now examine all quark contributions at orders $g$ and $g^{2}$ associated with a single-gluon
incoming state $\ket{\gl{a}{i}{k}}$. The possible final configurations are a quark--antiquark pair
(order $g$), the one gluon state itself (through the quark loop) and the quark--antiquark--gluon
state (order $g^{2}$).

\subsubsection{Quark--antiquark outgoing}

To order $g$ only one diagram contributes: the incoming gluon splits into a pair.
\begin{center}
\begin{tikzpicture}
\draw[gl] (0,0.85) -- (1.7,0.85); \node[left] at (0,0.85) {\scriptsize$\gl{a}{i}{k}$};
\draw[qk] (1.7,0.85) -- (3.6,1.6); \node[right] at (3.6,1.6) {\scriptsize$\qk{\beta}{\lambda_{1}}{l}$};
\draw[qk] (3.6,0.1) -- (1.7,0.85); \node[right] at (3.6,0.1) {\scriptsize$\aqk{\alpha}{\lambda_{2}}{m}$};
\draw[cut] (2.6,-0.2) -- (2.6,1.95);
\end{tikzpicture}
\end{center}
\begin{multline}
\ket{(\Psi^{a}_{i}(k))_{q\bar q}}=\int\! d^{3}l\,d^{3}m\;
\ket{\qk{\beta}{\lambda_{1}}{l}\aqk{\alpha}{\lambda_{2}}{m}}\;
\frac{\bra{\qk{\beta}{\lambda_{1}}{l}\aqk{\alpha}{\lambda_{2}}{m}}H_{gq\bar q}\ket{\gl{a}{i}{k}}}
{E_{g}(k)-E_{q\bar q}(l,m)}
\\
=\frac{g}{8\pi^{3/2}}\int\! d^{3}l\,d^{3}m\;\ket{q\bar q}\;
\frac{t^{a}_{\beta\alpha}\,\dthree(k-l-m)\,\Ph^{i}_{\lambda_{1}\lambda_{2}}(k;l,m)}
{\sqrt{k^{+}}\;\big(\eps_{k}-\eps_{l}-\eps_{m}\big)} .
\label{eq:Pg}
\end{multline}

\subsubsection{One gluon outgoing: the quark loop}

At order $g^{2}$ the incoming gluon can fluctuate into a pair and recombine. Since momentum is
conserved along the soft line, this contribution is proportional to $\dthree(k-k')$ and has
$E_{i}=E_{f}$: it is the quark loop contribution to the gluon self-energy, and it enters the
coefficient $B_{3}$ of Paper~I only through its momentum diagonal part. As in Paper~I, its effect
at $\order{g^{2}}$ is fixed by the normalization of the dressed state. The norm of the component
\eqref{eq:Pg} is evaluated with the variables $z=l^{+}/k^{+}$, $\bm\kappa=\bm l-z\bm k$ of
Eq.~\eqref{eq:kappa}, for which $d^{3}l=k^{+}dz\,d^{2}\kappa$ and, for fixed $i$ after averaging over the direction of
$\bm\kappa$,
\begin{equation}
\sum_{\lambda_{1}\lambda_{2}}\big|\Ph^{i}_{\lambda_{1}\lambda_{2}}(k;l,m)\big|^{2}
=\frac{2\,\bm\kappa^{2}\,\big[z^{2}+(1-z)^{2}\big]}{(k^{+})^{2}z^{2}(1-z)^{2}} .
\label{eq:PhiSq}
\end{equation}
All powers of $\bm\kappa$ and $k^{+}$ cancel against the energy denominator and one finds
\begin{align}
\big\langle(\Psi^{a'}_{i'}(k'))_{q\bar q}\big|(\Psi^{a}_{i}(k))_{q\bar q}\big\rangle
&=\delta^{aa'}\delta_{ii'}\dthree(k-k')\,\frac{g^{2}\Nf T_{F}}{8\pi^{3}}
\int_{\Lambda/k^{+}}^{1-\Lambda/k^{+}}\!\!dz\,\big[z^{2}+(1-z)^{2}\big]\int\!\frac{d^{2}\kappa}{\bm\kappa^{2}}
\nonumber\\
&=\delta^{aa'}\delta_{ii'}\dthree(k-k')\,\frac{g^{2}\Nf T_{F}}{8\pi^{3}}\,\frac{2}{3}\int\!\frac{d^{2}\kappa}{\bm\kappa^{2}},
\label{eq:qloop}
\end{align}
with $T_{F}=1/2$; the longitudinal integral is finite and the cutoff $\Lambda$ may be removed.
Combined with the gluon loop of Paper~I, Eq.~(4.83) there, the order $g^{2}$ norm of the
two-parton components of a soft gluon is
\begin{equation}
\frac{g^{2}}{8\pi^{3}}\left[2N_{c}\log\frac{k^{+}}{\Lambda}-\frac{11N_{c}-4T_{F}\Nf}{6}\right]
\int\!\frac{d^{2}\kappa}{\bm\kappa^{2}},
\label{eq:gnorm}
\end{equation}
in which the familiar combination $11N_{c}-4T_{F}\Nf$ of the one-loop $\beta$-function appears as
the coefficient of the ultraviolet divergence. The transverse integral is scaleless and vanishes
in dimensional regularization. Therefore, exactly as in Paper~I, the quark loop does not modify
the normalization $N$ at $\order{g^{2}}$, and the momentum diagonal part of $B_{3}$ is unchanged.

\subsubsection{Quark, antiquark and one gluon outgoing}

At order $g^{2}$ there are seven connected contributions to the
$\ket{\qk{\beta}{\lambda_{1}}{l}\aqk{\alpha}{\lambda_{2}}{m}\gl{d}{n}{r}}$ component. We denote the
overall energy denominator by
\begin{equation}
D_{f}\equiv E_{i}-E_{f}=\eps_{k}-\eps_{l}-\eps_{m}-\eps_{r} .
\end{equation}

\paragraph{(a) The gluon splits into a pair, then the quark emits a gluon.}
\begin{center}
\begin{tikzpicture}
\draw[gl] (0,0.9) -- (1.3,0.9); \node[left] at (0,0.9) {\scriptsize$\gl{a}{i}{k}$};
\draw[qk] (1.3,0.9) -- (2.4,1.5); \node[below right] at (1.55,1.08) {\scriptsize$p$};
\draw[qk] (2.4,1.5) -- (3.9,1.9); \node[right] at (3.9,1.9) {\scriptsize$\qk{\beta}{\lambda_{1}}{l}$};
\draw[gl] (2.4,1.5) -- (3.9,1.1); \node[right] at (3.9,1.1) {\scriptsize$\gl{d}{n}{r}$};
\draw[qk] (3.9,0.1) -- (1.3,0.9); \node[right] at (3.9,0.1) {\scriptsize$\aqk{\alpha}{\lambda_{2}}{m}$};
\draw[cut] (1.85,-0.25) -- (1.85,2.25); \draw[cut] (3.2,-0.25) -- (3.2,2.25);
\end{tikzpicture}
\end{center}
\begin{multline}
\ket{(\Psi^{a}_{i}(k))^{1}_{gq\bar q}}=\int\! d^{3}p\,d^{3}l\,d^{3}m\,d^{3}r\;
\ket{q\bar q\,g}\;
\\
\frac{\bra{\qk{\beta}{\lambda_{1}}{l}\gl{d}{n}{r}}H_{gq\bar q}\ket{\qk{\gamma}{\lambda}{p}}}{D_{f}}\;
\frac{\bra{\qk{\gamma}{\lambda}{p}\aqk{\alpha}{\lambda_{2}}{m}}H_{gq\bar q}\ket{\gl{a}{i}{k}}}
{\eps_{k}-\eps_{p}-\eps_{m}} .
\end{multline}
Using \eqref{eq:MEqqg} and \eqref{eq:MEqg}, closing the intermediate helicity sum with
$\sum_{\lambda}\chi_{\lambda}\chi^{\dg}_{\lambda}=\mathbf 1$, and $p=l+r$,
\begin{multline}
\ket{(\Psi^{a}_{i}(k))^{1}_{gq\bar q}}=\frac{g^{2}}{64\pi^{3}}\int\! d^{3}l\,d^{3}m\,d^{3}r\;
\ket{q\bar q\,g}\\
\times\,(t^{d}t^{a})_{\beta\alpha}\,\frac{\dthree(k-l-m-r)}{\sqrt{r^{+}k^{+}}}\,
\frac{\chi^{\dg}_{\lambda_{1}}\Gam^{n}(r;l,p)\,\Gam^{i}(k;p,m)\chi_{\lambda_{2}}}
{D_{f}\,\big(\eps_{k}-\eps_{p}-\eps_{m}\big)} .
\label{eq:B3a}
\end{multline}

\paragraph{(b) The gluon splits into a pair, then the antiquark emits a gluon.}
\begin{center}
\begin{tikzpicture}
\draw[gl] (0,1.6) -- (1.3,1.6); \node[left] at (0,1.6) {\scriptsize$\gl{a}{i}{k}$};
\draw[qk] (1.3,1.6) -- (3.9,2.1); \node[right] at (3.9,2.1) {\scriptsize$\qk{\beta}{\lambda_{1}}{l}$};
\draw[qk] (2.4,0.95) -- (1.3,1.6); \node[below] at (1.85,1.15) {\scriptsize$m'$};
\draw[qk] (3.9,0.5) -- (2.4,0.95); \node[right] at (3.9,0.5) {\scriptsize$\aqk{\alpha}{\lambda_{2}}{m}$};
\draw[gl] (2.4,0.95) -- (3.9,1.35); \node[right] at (3.9,1.35) {\scriptsize$\gl{d}{n}{r}$};
\draw[cut] (1.85,0.2) -- (1.85,2.45); \draw[cut] (3.2,0.2) -- (3.2,2.45);
\end{tikzpicture}
\end{center}
With $m'=m+r$, the antiquark vertex \eqref{eq:MEaqg} and the same steps,
\begin{multline}
\ket{(\Psi^{a}_{i}(k))^{2}_{gq\bar q}}=-\frac{g^{2}}{64\pi^{3}}\int\! d^{3}l\,d^{3}m\,d^{3}r\;
\ket{q\bar q\,g}\\
\times\,(t^{a}t^{d})_{\beta\alpha}\,\frac{\dthree(k-l-m-r)}{\sqrt{r^{+}k^{+}}}\,
\frac{\chi^{\dg}_{\lambda_{1}}\Gam^{i}(k;l,m')\,\Gam^{n}(r;m',m)\chi_{\lambda_{2}}}
{D_{f}\,\big(\eps_{k}-\eps_{l}-\eps_{m'}\big)} .
\label{eq:B3b}
\end{multline}

\paragraph{(c) The gluon splits into two gluons, one of which splits into a pair.}
\begin{center}
\begin{tikzpicture}
\draw[gl] (0,1.1) -- (1.3,1.1); \node[left] at (0,1.1) {\scriptsize$\gl{a}{i}{k}$};
\draw[gl] (1.3,1.1) -- (3.9,2.0); \node[right] at (3.9,2.0) {\scriptsize$\gl{d}{n}{r}$};
\draw[gl] (1.3,1.1) -- (2.4,0.6); \node[below] at (1.85,0.65) {\scriptsize$\gl{b'}{j'}{k'}$};
\draw[qk] (2.4,0.6) -- (3.9,1.15); \node[right] at (3.9,1.15) {\scriptsize$\qk{\beta}{\lambda_{1}}{l}$};
\draw[qk] (3.9,-0.05) -- (2.4,0.6); \node[right] at (3.9,-0.05) {\scriptsize$\aqk{\alpha}{\lambda_{2}}{m}$};
\draw[cut] (1.85,-0.4) -- (1.85,2.35); \draw[cut] (3.2,-0.4) -- (3.2,2.35);
\end{tikzpicture}
\end{center}
\begin{equation}
\ket{(\Psi^{a}_{i}(k))^{3}_{gq\bar q}}=\frac12\int\! d^{3}k'\,d^{3}l\,d^{3}m\,d^{3}r\;
\ket{q\bar q\,g}\;
\frac{\bra{q\bar q}H_{gq\bar q}\ket{\gl{b'}{j'}{k'}}}{D_{f}}\;
\frac{\bra{\gl{b'}{j'}{k'}\gl{d}{n}{r}}H_{ggg}\ket{\gl{a}{i}{k}}}{\eps_{k}-\eps_{k'}-\eps_{r}} .
\end{equation}
The factor $\tfrac12$ cancels against the two assignments of the intermediate gluon pair. With
$k'=l+m$ and \eqref{eq:MEggg},
\begin{multline}
\ket{(\Psi^{a}_{i}(k))^{3}_{gq\bar q}}=\frac{ig^{2}}{64\pi^{3}}\int\! d^{3}l\,d^{3}m\,d^{3}r\;
\ket{q\bar q\,g}\\
\times\,f^{ab'd}t^{b'}_{\beta\alpha}\,\frac{\dthree(k-l-m-r)}{k'^{+}\sqrt{k^{+}r^{+}}}\,
\frac{\Vg^{ij'n}(k;k',r)\;\Ph^{j'}_{\lambda_{1}\lambda_{2}}(k';l,m)}
{D_{f}\,\big(\eps_{k}-\eps_{k'}-\eps_{r}\big)} .
\label{eq:B3c}
\end{multline}

\paragraph{(d), (e) The valence emits the gluon $r$ and the incoming gluon splits.}
\begin{center}
\begin{tikzpicture}
\draw[val] (0,0) -- (3.9,0); \node[left] at (-0.05,0) {\scriptsize$\rho$};
\draw[gl] (0,1.7) -- (2.3,1.7); \node[left] at (0,1.7) {\scriptsize$\gl{a}{i}{k}$};
\draw[gl] (1.0,0) -- (3.9,0.9); \node[right] at (3.9,0.9) {\scriptsize$\gl{d}{n}{r}$};
\draw[qk] (2.3,1.7) -- (3.9,2.15); \node[right] at (3.9,2.15) {\scriptsize$\qk{\beta}{\lambda_{1}}{l}$};
\draw[qk] (3.9,1.4) -- (2.3,1.7); \node[right] at (3.9,1.4) {\scriptsize$\aqk{\alpha}{\lambda_{2}}{m}$};
\node at (1.9,-0.55) {\scriptsize (d)};
\end{tikzpicture}
\hspace{0.7cm}
\begin{tikzpicture}
\draw[val] (0,0) -- (3.9,0); \node[left] at (-0.05,0) {\scriptsize$\rho$};
\draw[gl] (0,1.7) -- (1.2,1.7); \node[left] at (0,1.7) {\scriptsize$\gl{a}{i}{k}$};
\draw[gl] (2.6,0) -- (3.9,0.9); \node[right] at (3.9,0.9) {\scriptsize$\gl{d}{n}{r}$};
\draw[qk] (1.2,1.7) -- (3.9,2.15); \node[right] at (3.9,2.15) {\scriptsize$\qk{\beta}{\lambda_{1}}{l}$};
\draw[qk] (3.9,1.4) -- (1.2,1.7); \node[right] at (3.9,1.4) {\scriptsize$\aqk{\alpha}{\lambda_{2}}{m}$};
\node at (1.9,-0.55) {\scriptsize (e)};
\end{tikzpicture}
\end{center}
In (d) the valence emits first and the intermediate state is
$\ket{\gl{a}{i}{k}\gl{d}{n}{r}}$, with $E_{i}-E_{m}=-\eps_{r}$; in (e) the pair is formed first, with
$E_{i}-E_{m}=\eps_{k}-\eps_{l}-\eps_{m}$:
\begin{align}
\ket{(\Psi^{a}_{i}(k))^{4}_{gq\bar q\,\rho}}&=\int\!d^{3}l\,d^{3}m\,d^{3}r\;\ket{q\bar q\,g}\;
\frac{\bra{q\bar q\,\gl{d}{n}{r}}H_{gq\bar q}\ket{\gl{a}{i}{k}\gl{d}{n}{r}}}{D_{f}}\;
\frac{\bra{\gl{a}{i}{k}\gl{d}{n}{r}}H_{g}\ket{\gl{a}{i}{k}}}{-\eps_{r}},
\\
\ket{(\Psi^{a}_{i}(k))^{5}_{gq\bar q\,\rho}}&=\int\!d^{3}l\,d^{3}m\,d^{3}r\;\ket{q\bar q\,g}\;
\frac{\bra{q\bar q\,\gl{d}{n}{r}}H_{g}\ket{q\bar q}}{D_{f}}\;
\frac{\bra{q\bar q}H_{gq\bar q}\ket{\gl{a}{i}{k}}}{\eps_{k}-\eps_{l}-\eps_{m}} .
\end{align}
Since $D_{f}=(\eps_{k}-\eps_{l}-\eps_{m})+(-\eps_{r})$, the two orderings combine by
\eqref{eq:factor} and the overall denominator cancels:
\begin{equation}
\ket{(\Psi^{a}_{i}(k))^{4+5}_{gq\bar q\,\rho}}=\frac{g^{2}}{32\pi^{3}}\int\!d^{3}l\,d^{3}m\,d^{3}r\;
\ket{q\bar q\,g}\;
\frac{t^{a}_{\beta\alpha}\,r^{n}\rho^{d}(-r)\,\Ph^{i}_{\lambda_{1}\lambda_{2}}(k;l,m)\,\dthree(k-l-m)}
{\sqrt{k^{+}}(r^{+})^{3/2}\,(-\eps_{r})\,\big(\eps_{k}-\eps_{l}-\eps_{m}\big)} .
\label{eq:B3de}
\end{equation}

\paragraph{(f) Instantaneous quark exchange.}
\begin{center}
\begin{tikzpicture}
\draw[gl] (0,1.9) -- (1.3,1.3); \node[left] at (0,1.9) {\scriptsize$\gl{a}{i}{k}$};
\draw[gl] (2.3,0.75) -- (3.9,0.15); \node[right] at (3.9,0.15) {\scriptsize$\gl{d}{n}{r}$};
\draw[qk] (1.3,1.3) -- (3.9,1.9); \node[right] at (3.9,1.9) {\scriptsize$\qk{\beta}{\lambda_{1}}{l}$};
\draw[ins] (1.3,1.3) -- (2.3,0.75); \ibarh{1.8}{1.025}
\draw[qk] (3.9,1.0) -- (2.3,0.75); \node[right] at (3.9,1.0) {\scriptsize$\aqk{\alpha}{\lambda_{2}}{m}$};
\draw[cut] (3.2,0.0) -- (3.2,2.2);
\end{tikzpicture}
\end{center}
The two orderings of the gluon attachments give instantaneous fermion lines with plus momenta
$l^{+}+r^{+}$ and $-(m^{+}+r^{+})$; using \eqref{eq:MEggqq1},
\begin{multline}
\ket{(\Psi^{a}_{i}(k))^{6}_{gq\bar q}}=\frac{g^{2}}{16\pi^{3}}\int\!d^{3}l\,d^{3}m\,d^{3}r\;
\ket{q\bar q\,g}\\
\times\,
\frac{\dthree(k-l-m-r)}{\sqrt{k^{+}r^{+}}\,D_{f}}\,
\chi^{\dg}_{\lambda_{1}}\!\left[\frac{(t^{d}t^{a})_{\beta\alpha}\sigma^{n}\sigma^{i}}{l^{+}+r^{+}}
-\frac{(t^{a}t^{d})_{\beta\alpha}\sigma^{i}\sigma^{n}}{m^{+}+r^{+}}\right]\chi_{\lambda_{2}} .
\label{eq:B3f}
\end{multline}

\paragraph{(g) Instantaneous gluon exchange between the gluon and quark currents.}
\begin{center}
\begin{tikzpicture}
\draw[gl] (0,1.6) -- (1.8,1.6); \node[left] at (0,1.6) {\scriptsize$\gl{a}{i}{k}$};
\draw[gl] (1.8,1.6) -- (3.9,2.05); \node[right] at (3.9,2.05) {\scriptsize$\gl{d}{n}{r}$};
\draw[ins] (1.8,1.6) -- (1.8,0.7); \ibar{1.8}{1.15}
\draw[qk] (1.8,0.7) -- (3.9,1.15); \node[right] at (3.9,1.15) {\scriptsize$\qk{\beta}{\lambda_{1}}{l}$};
\draw[qk] (3.9,0.1) -- (1.8,0.7); \node[right] at (3.9,0.1) {\scriptsize$\aqk{\alpha}{\lambda_{2}}{m}$};
\draw[cut] (3.2,-0.15) -- (3.2,2.35);
\end{tikzpicture}
\end{center}
This is the $J_{g}\cdot J_{q}$ term of \eqref{eq:Hqinst}; the exchanged gluon carries
$P^{+}=l^{+}+m^{+}=k^{+}-r^{+}$ and, with \eqref{eq:MEJJ1},
\begin{equation}
\ket{(\Psi^{a}_{i}(k))^{7}_{gq\bar q}}=\frac{ig^{2}}{2(2\pi)^{3}}\int\!d^{3}l\,d^{3}m\,d^{3}r\;
\ket{q\bar q\,g}\,
\frac{f^{bad}\,t^{b}_{\beta\alpha}\,\delta_{\lambda_{1}\lambda_{2}}\,\delta_{in}\,(k^{+}+r^{+})\,
\dthree(k-l-m-r)}{\sqrt{k^{+}r^{+}}\,(l^{+}+m^{+})^{2}\,D_{f}} .
\label{eq:B3g}
\end{equation}

The remaining topology, in which the valence emits a gluon which splits into the pair while the
incoming gluon is a spectator, is disconnected: it is the $\mf A_{1}$ term of \eqref{eq:Omg} and is
not part of $\mf B_{3}$.

\subsubsection{Summary of results}
\label{sec:sum-g}

The quark components of the dressed one gluon state through $\order{g^{2}}$ are
\begin{equation}
\ket{\Psi^{a}_{i}(k)}\supset\ket{(\Psi^{a}_{i}(k))_{q\bar q}}
+\sum_{n=1,2,3,6,7}\ket{(\Psi^{a}_{i}(k))^{n}_{gq\bar q}}
+\ket{(\Psi^{a}_{i}(k))^{4+5}_{gq\bar q\,\rho}}
+\ket{\gl{a}{i}{k}}\otimes\ket{\Psi_{q\bar q\,\rho}} ,
\label{eq:sum-g}
\end{equation}
where the last term is the vacuum pair component \eqref{eq:Pvac} with the incoming gluon as a
spectator. The diagrams (a)--(c) are the sequential processes $g\to q\bar q\to q\bar q\,g$ and
$g\to gg\to q\bar q\,g$, (d)+(e) is the product of the Weizs\"acker--Williams emission and the
splitting $g\to q\bar q$, and (f), (g) are the two instantaneous interactions. The overall
denominator $D_{f}$ cannot vanish except on the collinear boundary of phase space (a set of
measure zero), so $\mf B_{3}$ has no on-shell pole. As discussed above, the normalization is
unchanged at this order.
